\textbf{Phần 1. Trắc nghiệm lựa chọn}
\setcounterpageref{ex}{0}

\begin{ex} %[2D8V45]
Tính giá trị của biểu thức: $2 \times 8 + 35 : 5$.
\choice
{\True $23$}
{$21$}
{$22$}
{$24$}
\loigiai{
$2 \times 8 = 16$; $35 : 5 = 7$; $16 + 7 = 23$.
\textbf{Đáp án của câu hỏi là A}
}
\end{ex}

\begin{ex} %[2D8V45]
Có $18$ chiếc kẹo chia đều cho Mai và Lan. Sau đó mẹ cho Mai thêm $3$ cái kẹo nữa. Hỏi Mai có tất cả bao nhiêu cái kẹo?
\choice
{$10$ cái kẹo}
{$11$ cái kẹo}
{\True $12$ cái kẹo}
{$13$ cái kẹo}
\loigiai{
Mỗi bạn nhận được: $18 : 2 = 9$ (cái kẹo).\nMai có tất cả: $9 + 3 = 12$ (cái kẹo).
\textbf{Đáp án của câu hỏi là C}
}
\end{ex}

\begin{ex} %[2D8V45]
Tìm $x$ biết: $x : 2 = 15 - 8$.
\choice
{$12$}
{\True $14$}
{$16$}
{$18$}
\loigiai{
Vế phải: $15 - 8 = 7$. Ta có: $x = 7 \times 2 = 14$.
\textbf{Đáp án của câu hỏi là B}
}
\end{ex}

\begin{ex} %[2D8V45]
Một chuồng nuôi $5$ con thỏ và $6$ con gà. Hỏi cả chuồng có bao nhiêu cái chân?
\choice
{$30$ cái chân}
{$28$ cái chân}
{$34$ cái chân}
{\True $32$ cái chân}
\loigiai{
Số chân thỏ: $4 \times 5 = 20$ cái. Số chân gà: $2 \times 6 = 12$ cái. Tổng: $20 + 12 = 32$ cái chân.
\textbf{Đáp án của câu hỏi là D}
}
\end{ex}

\begin{ex} %[2D8V45]
Tính nhanh: $5 \times 4 + 5 \times 6$.
\choice
{\True $50$}
{$45$}
{$55$}
{$40$}
\loigiai{
$5 \times (4 + 6) = 5 \times 10 = 50$.
\textbf{Đáp án của câu hỏi là A}
}
\end{ex}

\begin{ex} %[2D8V45]
Có $40\text{ kg}$ gạo đóng đều vào các túi loại $5\text{ kg}$. Người ta đã bán đi $3$ túi. Hỏi còn lại bao nhiêu túi gạo?
\choice
{$4$ túi}
{$6$ túi}
{\True $5$ túi}
{$7$ túi}
\loigiai{
Tổng số túi: $40 : 5 = 8$ (túi). Còn lại: $8 - 3 = 5$ (túi).
\textbf{Đáp án của câu hỏi là C}
}
\end{ex}

\begin{ex} %[2D8V45]
Điền dấu thích hợp: $2 \times 9 \dots 5 \times 4 - 3$.
\choice
{$<$}
{\True $>$}
{$=$}
{Không so sánh được}
\loigiai{
Vế trái: $2 \times 9 = 18$. Vế phải: $20 - 3 = 17$. Vì $18 > 17$ nên điền dấu $>$. 
\textbf{Đáp án của câu hỏi là B}
}
\end{ex}

\begin{ex} %[2D8V45]
Một cuộn dây dài $20\text{ m}$ cắt thành các đoạn dài $2\text{ m}$. Hỏi phải cắt bao nhiêu nhát kéo?
\choice
{$10$ nhát}
{$8$ nhát}
{$7$ nhát}
{\True $9$ nhát}
\loigiai{
Số đoạn dây: $20 : 2 = 10$ đoạn. Số nhát cắt: $10 - 1 = 9$ nhát.
\textbf{Đáp án của câu hỏi là D}
}
\end{ex}

\begin{ex} %[2D8V45]
Tìm một số biết rằng lấy số đó nhân với $2$ rồi cộng với $14$ thì được $30$.
\choice
{\True $8$}
{$7$}
{$9$}
{$6$}
\loigiai{
Tính ngược: $30 - 14 = 16$; $16 : 2 = 8$.
\textbf{Đáp án của câu hỏi là A}
}
\end{ex}

\begin{ex} %[2D8V45]
Trong một phép chia cho $5$, thương là số chẵn lớn nhất có một chữ số. Số bị chia là
\choice
{$35$}
{$45$}
{\True $40$}
{$30$}
\loigiai{
Số chẵn lớn nhất có 1 chữ số là 8. Số bị chia: $8 \times 5 = 40$.
\textbf{Đáp án của câu hỏi là C}
}
\end{ex}

\begin{ex} %[2D8V45]
Mỗi can chứa $5\text{ L}$ dầu. Bác Nam có $35\text{ L}$ dầu đóng vào các can như thế và cho bạn $2$ can. Bác còn lại bao nhiêu can dầu?
\choice
{$4$ can}
{\True $5$ can}
{$6$ can}
{$7$ can}
\loigiai{
Bác đóng được: $35 : 5 = 7$ (can). Còn lại: $7 - 2 = 5$ (can).
\textbf{Đáp án của câu hỏi là B}
}
\end{ex}

\begin{ex} %[2D8V45]
Hai số có tích bằng $18$ và hiệu bằng $7$. Hai số đó là
\choice
{$6\text{ và } 3$}
{$18\text{ và } 1$}
{$8\text{ và } 1$}
{\True $9\text{ và } 2$}
\loigiai{
$9 \times 2 = 18$ và $9 - 2 = 7$. Hai số là 9 và 2.
\textbf{Đáp án của câu hỏi là D}
}
\end{ex}

\textbf{Phần 2. Trắc nghiệm đúng sai}
\setcounterpageref{ex}{0}

\begin{ex} %[2D8V45]
Xét tính đúng sai của các bài toán rèn luyện tư duy nhân, chia:
\choiceTF
{\True Biểu thức 2 x 7 + 5 x 4 có kết quả bằng 34.}
{Thực hiện phép tính 40 : 5 - 3 ta lấy 5 - 3 trước rồi lấy 40 chia.}
{\True Số chia luôn khác 0.}
{Tích của 2 số tự nhiên luôn lớn hơn thương của 2 số đó.}
\loigiai{
\begin{itemize}
    \item Ý a) Đúng: $14 + 20 = 34$. Đúng.
    \item Ý b) Sai: Quy tắc: Nhân chia trước, cộng trừ sau. Phải lấy $40 : 5 = 8$ rồi $8 - 3 = 5$.
    \item Ý c) Đúng: Đúng.
    \item Ý d) Sai: Nếu $2 \times 1 = 2$ và $2 : 1 = 2$, tích bằng thương.
\end{itemize}
\begin{itemize}
    \item a) \textbf{Đúng}.
    \item b) \textbf{Sai}.
    \item c) \textbf{Đúng}.
    \item d) \textbf{Sai}.
\end{itemize}
}
\end{ex}

\begin{ex} %[2D8V45]
Bác thợ mộc có 16 m gỗ, cưa thành các đoạn dài 2 m:
\choiceTF
{\True Bác cưa được 8 đoạn gỗ.}
{\True Bác phải cưa 7 nhát cưa.}
{Nếu mỗi nhát cưa mất 2 phút thì bác cưa hết 16 phút.}
{\True Nếu mỗi đoạn gỗ bác đóng thành một chiếc ghế thì bác đóng được 8 chiếc ghế.}
\loigiai{
\begin{itemize}
    \item Ý a) Đúng: $16 : 2 = 8$ đoạn.
    \item Ý b) Đúng: $8 - 1 = 7$ nhát.
    \item Ý c) Sai: Thời gian cưa: $7 \times 2 = 14$ phút (không phải 16 phút).
    \item Ý d) Đúng: Đúng.
\end{itemize}
\begin{itemize}
    \item a) \textbf{Đúng}.
    \item b) \textbf{Đúng}.
    \item c) \textbf{Sai}.
    \item d) \textbf{Đúng}.
\end{itemize}
}
\end{ex}

\textbf{Phần 3. Trả lời ngắn}
\setcounterpageref{ex}{0}

\begin{ex} %[2D8V45]
Tính giá trị biểu thức: 2 x 9 + 45 : 5.
\shortans{27}
\loigiai{
\textbf{Đáp số:} $27$\\
$2 \times 9 = 18$; $45 : 5 = 9$; $18 + 9 = 27$.
}
\end{ex}

\begin{ex} %[2D8V45]
Có 5 bạn nhỏ, mỗi bạn có 4 quả bóng bay. Sau đó bị nổ mất 6 quả. Hỏi còn lại bao nhiêu quả bóng bay?
\shortans{14}
\loigiai{
\textbf{Đáp số:} $14$\\
$5 \times 4 - 6 = 20 - 6 = 14$ quả.
}
\end{ex}

\begin{ex} %[2D8V45]
Tìm x biết: x : 5 + 17 = 23.
\shortans{30}
\loigiai{
\textbf{Đáp số:} $30$\\
$x : 5 = 23 - 17 = 6 \Rightarrow x = 6 \times 5 = 30$.
}
\end{ex}

\begin{ex} %[2D8V45]
Một chuồng nuôi gà có tất cả 16 cái chân. Hỏi chuồng đó có bao nhiêu con gà?
\shortans{8}
\loigiai{
\textbf{Đáp số:} $8$\\
$16 : 2 = 8$ con gà.
}
\end{ex}

\begin{ex} %[2D8V45]
Bao gạo nặng 50 kg chia đều vào 5 túi. Mẹ đã dùng hết 2 túi. Hỏi còn lại bao nhiêu ki-lô-gam gạo?
\shortans{30}
\loigiai{
\textbf{Đáp số:} $30$\\
Mỗi túi: $50 : 5 = 10\text{ kg}$. Còn 3 túi: $10 \times 3 = 30\text{ kg}$.
}
\end{ex}

\begin{ex} %[2D8V45]
Tìm một số biết rằng lấy số đó chia cho 2 rồi nhân với 5 thì được 20.
\shortans{8}
\loigiai{
\textbf{Đáp số:} $8$\\
$20 : 5 = 4$; $4 \times 2 = 8$.
}
\end{ex}
